\scanpage{46}
Soit $\Sigma$ un ens. de segments de $S$, satisfaisant la condition suivante

($\deleted{\unsure{\Sigma|>}}\sigma$ 1) Si $A,B\in\Sigma$, et si $A\cap B\ne\phi$, alors $A\subset B$ ou $B\subset A$. \inserted{($\sigma.2$)} Si de plus $A\ne B$ et $A\subset B$, alors $A$ ne contient pas d'extrémité de $B$.

Posons alors, pour $A\in\Sigma$
\[
 A^*=\unsure{.|\cdot}\, A\setminus\bigcup_{\substack{B\in\Sigma\\ B\subset A\\ B\ne A}} B
\]
On a (par $\sigma.2$)
\[
 A^*\supset\partial A\quad\text{(ens des extrémités de $A$)}
\]
donc $A^*\ne\phi$. Je dis que
\[
 \begin{aligned}
 A&\longmapsto A^*\\
 \Sigma&\longrightarrow \unsure{P|\mathcal P|\mathfrak P}(S)
 \end{aligned}
 \qquad\text{est injectif}
\]
% slanted
\margin{\uline{Prop} Si $A,B\in\Sigma$, $A\ne B$, alors $A^*\cap B^*=\phi$ (donc $A^*\ne B^*$). \deleted{\unsure{sont disj.|sont di}} et $A^*,B^*$ non panachés. On a $\bigcup_{A\in\Sigma} A=\bigcup_{A\in\Sigma} A^*$. Si $x\in\bigcup_{A\in\Sigma}A$, alors l'ens \inserted{$\Sigma(x)$} des $A\in\Sigma$ tels que $x\in\Sigma$ est tot. ordonné, donc admet un plus petit \unsure{élément|élém\textsuperscript{t}}, soit $A$, et alors $A^*\ni x$.}
i.e. que si $A,B\in\Sigma$, alors
\[
 A^*=B^*\quad\Longrightarrow\quad A=B
\]
\inserted{i.e. $A\ne B$ \deleted{distincts} $\Longrightarrow A^*\ne B^*$. Comme $A^*,B^*$ sont $\ne\phi$, il suffit de prouver $A^*\cap B^*=\phi$.} \deleted{En effet}, si $A\cap B=\phi$, alors $A^*\cap B^*=\phi$, donc $A^*\ne B^*$ (puisque $A^*,B^*$ non vides). Si $A\cap B\ne\phi$, \unsure{on a|$\cdots$} \inserted{(par $\sigma\,1$)} $A\subset B$ ou $B\subset A$, soit $A\subset B$.

\deleted{Mais alors $A\cap\partial B=\phi$ \inserted{par $\sigma 2$}, \deleted{donc} $A^*\cap\partial B=\phi$, donc \unsure{$A^*\supset$|$A\supset$} (comme $\phi\ne\partial B\subset B^*$) $A^*\not\supset B^*$. Comme \unsure{on aurait|on a} $A^*\cap\partial B=$} Comme $A\ne B$, $A\cap B^*=\phi$ par définition de $B^*$, donc $A^*\cap B^*=\phi$, \unsure{qfd.|cqfd.} \deleted{Je dis} Aussi, $A^*,B^*$ non panachés, puisque $A^*\subset A$ segment $\subset S\setminus B^*$.
% archive sheet number: 45
